\cotefile{10}
\scanpage{102}
% Archive pencil foliation: 101.

Quelques exemples de catégories tensorielles.

1) Soit $M$ un groupe. Soit $\underline{C}_M$\rlap{\handwritten{\tikz[baseline=0pt,x=1em,y=1em]{\draw (0,1.25)--(.4,-.65);}}} la catégorie des vectoriels de dimension finie sur le corps $k$, munis d'une graduation de type $M$. C'est une catégorie tensorielle sur $k$, munie d'un foncteur fibre sur $k$, le foncteur ``oublif de la graduation''. Le groupe algébrique associé est le groupe de type multiplicatif $D_k(M)$ (SGA 3 I 4.7.3). Par exemple si $M=\underline{Z}^r$, on trouve $G=\underline{G}_m^r$.

Application: soit $\underline{C}$ une catégorie tensorielle sur $k$ munie d'un foncteur fibre \rlap{\handwritten{\tikz[baseline=-.1ex,x=1em,y=1em]{\draw (-.55,.7) .. controls (-.25,.35) .. (0,-.25) .. controls (.2,.35) .. (.65,1.1);}}}\inserted{$F$} sur $k$, donc associée à un schéma en groupes affine $G$ sur $k$\unsure{,|.} On cherche toutes les façons de mettre, pour chaque $\deleted{M}\inserted{\handwritten{V}}\handwritten{\in}\operatorname{Ob}\underline{C}$, une graduation de type $M$ sur $F(V)$, de façon fonctorielle en $M$, et compatible (dans un sens évident) ave\unsure{\deleted{u}c|\deleted{c}u} les produits tensoriels.El\unsure{l|I}\unsure{\deleted{l}e|\deleted{e}l}s correspondent aux $\otimes$-foncteurs de $\underline{C}$ dans $\underline{C}_M$ compatibles avec les foncteurs fibres, donc aux homomorphismes de $D_k(M)$ dans $G$. Par exemple, si $M=\underline{Z}^r$, il faut prendre les homomorphismes de $\underline{\unsure{\deleted{Z}G|\deleted{G}Z}}_m^r$ dans $G$.

Dans la situation précédente, on peut se demander quand une $\otimes$-graduation de type $M$ du foncteur $F$ correspond à une graduation de type $M$ du foncteur identique de $\underline{C}$, i.e. pour tout $V$, la graduation de $F(V)$ provient d'une graduation de $V$. On trouve qu'il faut et il suffit pour cela que l'homomorphisme correspondant \begin{tikzcd}[cramped]D_k(M)\arrow[r]&G\end{tikzcd} soit central. Par exe\unsure{m|n}ple, si $\underline{C}$ est la catégorie des m\unsure{\deleted{t}o|\deleted{o}t}\unsure{\deleted{o}t|\deleted{t}o}\unsure{\deleted{f}i|\deleted{i}f}fs sur $k$, et si on dispose d'un foncteur fibre $F$ de $\underline{C}$ sur $k$, alors on trouve un homomorphisme central canonique $i:$\begin{tikzcd}[cramped]\underline{G}_m\arrow[r]&G\end{tikzcd}. D'ailleurs, la donnée du motif de Tate (qui est de rang $2$, èt de ``poids'' $2$ pour la graduation naturelle) correspond à \deleted{un} la donnée d'un homomorphisme $j:$\begin{tikzcd}[cramped]G\arrow[r]&\underline{G}_m=\operatorname{Gl}(1)\end{tikzcd}. Le fait que $T$ soit de poids $2$ s'exprime par la relation
\[
ji(\handwritten{\lambda})=\handwritten{\lambda}^{2}.
\]

